% =====================================================================
%  structure.tex  —  Преамбула учебника по ИИ для 10–11 классов
%
%  Совместимость: документ собирается БЕЗ изменений и под pdfLaTeX,
%  и под XeLaTeX/LuaLaTeX. Компилятор определяется автоматически
%  через пакет iftex.
%
%  Возможности: кириллица, продвинутая математика, рисунки TikZ/PGF,
%  гиперссылки, листинги кода Python, оформленные «врезки»
%  (теоремы, определения, исторические справки).
% =====================================================================

% ---------------------------------------------------------------------
%  ВЫБОР КОМПИЛЯТОРА
% ---------------------------------------------------------------------
\usepackage{iftex}

\ifPDFTeX
% ====================  pdfLaTeX  ====================================
%
% Классический путь: cm-super (Latin Modern в T2A-кодировке) +
% babel + UTF-8 на входе. Гарантированно работает на любой
% полной поставке TeX Live, включая Overleaf, без выбора компилятора.
%
\usepackage{cmap}                         % поисковый текст в PDF
\usepackage[utf8]{inputenc}
\usepackage[T2A,T1]{fontenc}
\usepackage[english,main=russian]{babel}
\usepackage{lmodern}

\else
% ====================  XeLaTeX / LuaLaTeX  ==========================
%
% Современный путь: polyglossia + fontspec с системными Unicode-шрифтами
% и поддержкой OpenType. Даёт более качественную типографику и
% свободу выбора шрифта.
%
\usepackage{polyglossia}
\setdefaultlanguage{russian}
\setotherlanguage{english}

\usepackage{fontspec}
\defaultfontfeatures{Ligatures=TeX}

% Fallback-цепочка шрифтов: пробуем найти первый доступный.
% На Overleaf ловится CMU Serif (пакет cm-unicode), на локальных
% Linux — обычно Liberation Serif. Если ни одного из перечисленных
% нет, остаётся Latin Modern (без кириллицы — будет видно).
\IfFontExistsTF{CMU Serif}{%
    \setmainfont{CMU Serif}%
}{\IfFontExistsTF{Libertinus Serif}{%
    \setmainfont{Libertinus Serif}%
}{\IfFontExistsTF{Linux Libertine O}{%
    \setmainfont{Linux Libertine O}%
}{\IfFontExistsTF{PT Serif}{%
    \setmainfont{PT Serif}%
}{\IfFontExistsTF{Liberation Serif}{%
    \setmainfont{Liberation Serif}%
}{%
    \setmainfont{Latin Modern Roman}%
}}}}}

\IfFontExistsTF{CMU Sans Serif}{%
    \setsansfont{CMU Sans Serif}%
}{\IfFontExistsTF{Libertinus Sans}{%
    \setsansfont{Libertinus Sans}%
}{\IfFontExistsTF{PT Sans}{%
    \setsansfont{PT Sans}%
}{\IfFontExistsTF{Liberation Sans}{%
    \setsansfont{Liberation Sans}%
}{%
    \setsansfont{Latin Modern Sans}%
}}}}

\IfFontExistsTF{CMU Typewriter Text}{%
    \setmonofont{CMU Typewriter Text}[Scale=0.95]%
}{\IfFontExistsTF{Libertinus Mono}{%
    \setmonofont{Libertinus Mono}[Scale=0.95]%
}{\IfFontExistsTF{PT Mono}{%
    \setmonofont{PT Mono}[Scale=0.95]%
}{\IfFontExistsTF{Liberation Mono}{%
    \setmonofont{Liberation Mono}[Scale=0.92]%
}{\IfFontExistsTF{DejaVu Sans Mono}{%
    \setmonofont{DejaVu Sans Mono}[Scale=0.92]%
}{%
    \setmonofont{Latin Modern Mono}[Scale=0.95]%
}}}}}
\fi
% ====================================================================

% --- Геометрия страницы ----------------------------------------------
\usepackage[a4paper,
            inner=2.6cm, outer=2.4cm,
            top=2.2cm,   bottom=2.4cm,
            headheight=15pt]{geometry}

% --- Базовая математика ----------------------------------------------
\usepackage{amsmath, amssymb, amsthm, mathtools}
\usepackage{bm}                  % полужирные математические символы
\allowdisplaybreaks              % разрыв длинных формул через страницы

% --- Цвета, графика, рисунки -----------------------------------------
\usepackage{xcolor}
\definecolor{accent}{HTML}{1F4E79}     % тёмно-синий — основной акцент
\definecolor{accent2}{HTML}{B85450}    % тёплый красный — выделение
\definecolor{shade}{HTML}{F4F1EA}      % светлый фон врезок
\definecolor{codebg}{HTML}{F7F7F2}     % фон листингов
\definecolor{codekw}{HTML}{1F4E79}     % ключевые слова
\definecolor{codecm}{HTML}{6A737D}     % комментарии
\definecolor{codestr}{HTML}{B85450}    % строки

\usepackage{graphicx}
\graphicspath{{figures/}}

\usepackage{tikz}
\usetikzlibrary{calc, arrows.meta, decorations.pathreplacing,
                positioning, shapes.geometric}
\usepackage{pgfplots}
\pgfplotsset{compat=1.18}

\usepackage{caption}
\captionsetup{labelfont={bf,color=accent}, font=small,
              labelsep=period, format=plain}

% --- Колонтитулы и заголовки -----------------------------------------
\usepackage{fancyhdr}
\pagestyle{fancy}
\fancyhf{}
\fancyhead[LE]{\small\itshape\nouppercase{\leftmark}}
\fancyhead[RO]{\small\itshape\nouppercase{\rightmark}}
\fancyfoot[LE,RO]{\thepage}
\renewcommand{\headrulewidth}{0.4pt}
\renewcommand{\footrulewidth}{0pt}

\usepackage{titlesec}
\titleformat{\chapter}[display]
  {\normalfont\sffamily\Huge\bfseries\color{accent}}
  {\filright\Large Глава~\thechapter}
  {6pt}{\filright\Huge}
\titlespacing*{\chapter}{0pt}{-10pt}{20pt}

\titleformat{\section}
  {\normalfont\sffamily\Large\bfseries\color{accent}}
  {\thesection}{0.7em}{}
\titleformat{\subsection}
  {\normalfont\sffamily\large\bfseries\color{accent!85!black}}
  {\thesubsection}{0.6em}{}

% --- Гиперссылки -----------------------------------------------------
% Под pdfLaTeX опция unicode=true нужна, чтобы кириллица в закладках
% PDF и метаданных отображалась корректно. Под XeLaTeX это делается
% автоматически.
\ifPDFTeX
  \usepackage[unicode=true, pdfusetitle, colorlinks=true,
              linkcolor=accent, citecolor=accent,
              urlcolor=accent2, filecolor=accent]{hyperref}
\else
  \usepackage[pdfusetitle, colorlinks=true,
              linkcolor=accent, citecolor=accent,
              urlcolor=accent2, filecolor=accent]{hyperref}
\fi
\usepackage{bookmark}

\usepackage[noabbrev]{cleveref}
\crefname{equation}{формула}{формулы}
\Crefname{equation}{Формула}{Формулы}
\crefname{figure}{рис.}{рис.}
\Crefname{figure}{Рис.}{Рис.}
\crefname{chapter}{глава}{главы}
\Crefname{chapter}{Глава}{Главы}
\crefname{section}{раздел}{разделы}
\Crefname{section}{Раздел}{Разделы}
\crefname{subsection}{раздел}{разделы}
\Crefname{subsection}{Раздел}{Разделы}
\crefname{table}{табл.}{табл.}
\Crefname{table}{Табл.}{Табл.}
% tcolorbox-теоремы используют свой счётчик tcb@cnt@theorem.
\crefname{tcb@cnt@theorem}{теорема}{теоремы}
\Crefname{tcb@cnt@theorem}{Теорема}{Теоремы}

% --- Окружения «врезок» (теоремы, определения, замечания, примеры) ---
\usepackage[most]{tcolorbox}
\tcbuselibrary{theorems, breakable, skins}

\newtcbtheorem[number within=chapter]{theorem}{Теорема}
  {enhanced, breakable,
   colback=white, colframe=accent,
   fonttitle=\bfseries\sffamily, coltitle=white,
   left=3mm, right=3mm, top=2mm, bottom=2mm,
   sharp corners, boxrule=0.6pt,
   attach boxed title to top left={xshift=4mm, yshift=-3mm},
   boxed title style={colback=accent, sharp corners, boxrule=0pt},
   separator sign=. ,
   description delimiters={(}{)}}{thm}

\newtcbtheorem[use counter from=theorem]{definition}{Определение}
  {enhanced, breakable,
   colback=shade, colframe=accent,
   fonttitle=\bfseries\sffamily, coltitle=white,
   left=3mm, right=3mm, top=2mm, bottom=2mm,
   sharp corners, boxrule=0.4pt,
   attach boxed title to top left={xshift=4mm, yshift=-3mm},
   boxed title style={colback=accent, sharp corners, boxrule=0pt},
   separator sign=. }{def}

\newtcbtheorem[use counter from=theorem]{example}{Пример}
  {enhanced, breakable,
   colback=white, colframe=accent2,
   fonttitle=\bfseries\sffamily, coltitle=white,
   left=3mm, right=3mm, top=2mm, bottom=2mm,
   sharp corners, boxrule=0.5pt,
   attach boxed title to top left={xshift=4mm, yshift=-3mm},
   boxed title style={colback=accent2, sharp corners, boxrule=0pt},
   separator sign=. }{ex}

% «Историческая справка» — мягкая врезка
\newtcolorbox{historybox}[1][]{enhanced, breakable,
   colback=shade, colframe=shade,
   leftrule=3pt, left=3mm, right=3mm, top=2mm, bottom=2mm,
   sharp corners, boxrule=0pt, frame hidden,
   borderline west={3pt}{0pt}{accent2},
   fonttitle=\bfseries\sffamily, title={Историческая справка}, #1}

% «На полях» / замечание
\newtcolorbox{notebox}[1][]{enhanced, breakable,
   colback=white, colframe=accent!60!white,
   leftrule=3pt, left=3mm, right=3mm, top=2mm, bottom=2mm,
   sharp corners, boxrule=0pt, frame hidden,
   borderline west={3pt}{0pt}{accent!60!white},
   fonttitle=\bfseries\sffamily, title={Замечание}, #1}

% Окружение для доказательств — стандартное amsthm
\renewcommand{\proofname}{Доказательство}

% --- Листинги кода (Python) ------------------------------------------
% Под pdfLaTeX пакет listings не умеет работать с UTF-8 кириллицей
% (это известная неудобная особенность). Самое надёжное переносимое
% решение — писать комментарии в коде по-английски. Под XeLaTeX/LuaLaTeX
% кириллица работала бы, но мы используем единый стиль ради совместимости.
\usepackage{listings}
\lstdefinestyle{pythonstyle}{
   language=Python,
   basicstyle=\ttfamily\small,
   keywordstyle=\color{codekw}\bfseries,
   commentstyle=\color{codecm}\itshape,
   stringstyle=\color{codestr},
   numberstyle=\tiny\color{codecm},
   numbers=left, numbersep=8pt,
   backgroundcolor=\color{codebg},
   showstringspaces=false,
   breaklines=true, breakatwhitespace=true,
   frame=leftline, framerule=1.2pt, rulecolor=\color{accent},
   xleftmargin=12pt, framexleftmargin=8pt,
   tabsize=4,
   columns=flexible, keepspaces=true,
}
\lstset{style=pythonstyle}

% --- Удобные сокращения для математики -------------------------------
\newcommand{\R}{\mathbb{R}}
\newcommand{\N}{\mathbb{N}}
\newcommand{\Z}{\mathbb{Z}}
\newcommand{\eps}{\varepsilon}
\newcommand{\norm}[1]{\left\|#1\right\|}
\newcommand{\abs}[1]{\left|#1\right|}
\DeclareMathOperator*{\argmin}{arg\,min}
\DeclareMathOperator*{\argmax}{arg\,max}

% --- Список задач в конце параграфа ----------------------------------
\newtcolorbox{tasksbox}{enhanced, breakable,
   colback=white, colframe=accent!70!black,
   fonttitle=\bfseries\sffamily, title={Задачи для самостоятельной работы},
   left=3mm, right=3mm, top=2mm, bottom=2mm,
   sharp corners, boxrule=0.6pt}

% --- Микротипографика ------------------------------------------------
\usepackage{microtype}

% --- Метаданные PDF --------------------------------------------------
\hypersetup{
  pdftitle  = {Искусственный интеллект. Учебник для 10-11 классов},
  pdfauthor = {Авторский коллектив},
  pdfsubject= {Информатика, продвинутый уровень},
  pdfkeywords = {ИИ, машинное обучение, оптимизация, метод Ньютона}
}
